%% ma: L11-B16-0021
%% lop: 11
%% chuong: 5
%% bai: 16
%% dang: 11.16.5
%% loai: TN4
%% muc: TH
%% dap_an: "A"
%% nguon: "tai-lieu-hoc-tap-mon-toan-lop-11-theo-chuong-trinh-giao-duc-pho-thong-2018.pdf (Lớp Toán Thầy Duy, Phạm Lê Duy, THPT Chu Văn An), Chương 3 Bài 2, Phần 1, Câu 9 (tr.360)"
%% kien_thuc: "Lớp 11 Bài 16 (giới hạn của hàm số)"
%% trang_thai: chinh-thuc
%% ghi_chu: "Đã duyệt 03/11/2026."
\begin{ex}
Trong các giới hạn dưới đây, giới hạn nào là $+\infty$?
\choice{\True $\lim\limits_{x\to4^-}\dfrac{2x-1}{4-x}$.}{$\lim\limits_{x\to+\infty}\left(-x^3+2x+3\right)$.}{$\lim\limits_{x\to-\infty}\dfrac{x^2+x+1}{x-1}$.}{$\lim\limits_{x\to4^+}\dfrac{2x-1}{4-x}$.}
\loigiai{\sloppy A: tử $\to7>0$, mẫu $4-x\to0^+$ nên $+\infty$. B: $-\infty$. C: $\sim x\to-\infty$. D: tử $\to7>0$, mẫu $\to0^-$ nên $-\infty$.}
\end{ex}
